\subsection{代数的森田等价}

命题 13. —— a) 若两个 k-代数是森田等价的，则它们的中心是同构的 k-代数。

\begin{enumerate}
\def\labelenumi{\alph{enumi})}
\setcounter{enumi}{1}
\item
  两个交换 \(k\)-代数森田等价当且仅当它们同构。
\item
  两个都是体的 k-代数森田等价当且仅当它们同构。
\item
  对 \(i = 1,2\)，设 \(\mathrm{A}_i\) 与 \(\mathrm{B}_i\) 是森田等价的 \(k\)-代数，\(\mathrm{P}_i\) 是一个可逆 \((\mathrm{A}_i,\mathrm{B}_i)_k\)-双模。令 \(\mathrm{A} = \mathrm{A}_1\otimes_k\mathrm{A}_2\)，\(\mathrm{B} = \mathrm{B}_1\otimes_k\mathrm{B}_2\)，\(\mathrm{P} = \mathrm{P}_1 \otimes_k \mathrm{P}_2\)。\(k\)-代数 A 与 B 是森田等价的，且 P 是一个可逆 \((\mathrm{A}, \mathrm{B})_k\)-双模。
\item
  若 A 与 B 是森田等价的 \(k\)-代数，且 \(k'\) 是一个交换 \(k\)-代数，则 \(k'\)-代数 \(\mathrm{A}_{(k')}\) 与 \(\mathrm{B}_{(k')}\) 是森田等价的。
\end{enumerate}

断言 a) 由 VIII, 第 102 页的推论 1 得出，并蕴含 b)。

设 \(\mathrm{K}\) 与 \(\mathrm{L}\) 是都是体的 \(k\)-代数，\(\mathrm{P}\) 是一个可逆 \((\mathrm{K},\mathrm{L})_k\)-双模。右 \(\mathrm{L}\)-向量空间 \(\mathrm{P}\) 是一个单模（VIII, 第 104 页），因而是 1 维的，于是 \(k\)-代数 \(\operatorname{End}_{\mathrm{L}}(\mathrm{P})\) 与 \(\mathrm{L}\) 同构。由 VIII, 第 101 页的定理 1，映射 \(a\mapsto a_{\mathrm{P}}\)（从 \(\mathrm{K}\) 到 \(\operatorname{End}_{\mathrm{L}}(\mathrm{P})\)）是一个同构。因此，体 \(\mathrm{K}\) 与 \(\mathrm{L}\) 在 \(k\) 上同构；断言 c) 得证。

在 d) 的假设下，设 \(Q_{i}\)（i = 1, 2）是一个 \((\mathrm{B}_{i}, \mathrm{A}_{i})_{k}\)-双模，与 \(P_{i}\) 互逆。用 Q 表示 \((\mathrm{B}, \mathrm{A})_{k}\)-双模 \(Q_{1} \otimes_{k} Q_{2}\)。考虑典范 k-线性同构 \((\mathrm{P}_{1} \otimes_{k} \mathrm{P}_{2}) \otimes_{k} (\mathrm{Q}_{1} \otimes_{k} \mathrm{Q}_{2}) \to (\mathrm{P}_{1} \otimes_{k} \mathrm{Q}_{1}) \otimes_{k} (\mathrm{P}_{2} \otimes_{k} \mathrm{Q}_{2})\)；在过渡到商时，它定义一个态射

\[
\left(\mathrm{P} _ {1} \otimes_ {k} \mathrm{P} _ {2}\right) \otimes_ {\mathrm{B}} \left(\mathrm{Q} _ {1} \otimes_ {k} \mathrm{Q} _ {2}\right)\rightarrow \left(\mathrm{P} _ {1} \otimes_ {\mathrm{B} _ {1}} \mathrm{Q} _ {1}\right) \otimes_ {k} \left(\mathrm{P} _ {2} \otimes_ {\mathrm{B} _ {2}} \mathrm{Q} _ {2}\right)
\]

它是 (A,A)-线性的。反之，逆同构 \((\mathrm{P}_{1}\otimes_{k}\mathrm{Q}_{1})\otimes_{k}(\mathrm{P}_{2}\otimes_{k}\mathrm{Q}_{2})\to(\mathrm{P}_{1}\otimes_{k}\mathrm{P}_{2})\otimes_{k}(\mathrm{Q}_{1}\otimes_{k}\mathrm{Q}_{2})\) 定义一个 (A,A)-线性态射 \((\mathrm{P}_{1}\otimes_{\mathrm{B}_{1}}\mathrm{Q}_{1})\otimes_{k}(\mathrm{P}_{2}\otimes_{\mathrm{B}_{2}}\mathrm{Q}_{2})\to(\mathrm{P}_{1}\otimes_{k}\mathrm{P}_{2})\otimes_{\mathrm{B}}(\mathrm{Q}_{1}\otimes_{k}\mathrm{Q}_{2})\)。这两个态射互为逆映射，从而都是同构。由于 \((\mathrm{A}_{i},\mathrm{A}_{i})\)-双模 \(P_{i}\otimes_{B_{i}}Q_{i}\) 同构于 \(A_{i}\)，我们得到一个 (A,A)-线性同构 \(P\otimes_{B}Q\to A\)。我们同样得到一个 (B,B)-线性同构 \(Q\otimes_{A}P\to B\)，这就完成了 d) 的证明。

在 e) 的假设下，设 P 是一个可逆 \((\mathrm{A},\mathrm{B})_{k}\)-双模；则 \(\mathrm{P}_{(k^{\prime})}\) 是一个可逆 \((\mathrm{A}_{(k^{\prime})},\mathrm{B}_{(k^{\prime})})_{k^{\prime}}\)-双模。

设 A 是一个 k-代数，e 是 A 中的一个幂等元。集合 eAe 赋以由 A 的加法、乘法与 k 的作用所诱导的运算，是一个以 e 为单位元的 k-代数。

命题 14. —— 设 A 与 B 是 k-代数。那么 A 与 B 森田等价当且仅当存在一个整数 \(n \geqslant 1\) 与一个方阵 \(e = (e_{ij})\)（在 \(\mathbf{M}_{n}(\mathbf{B})\) 中），满足下列条件：

\begin{enumerate}
\def\labelenumi{(\roman{enumi})}
\item
  我们有 \(e^2 = e\)。
\item
  由诸元素 \(e_{ij}\) 生成的 B 的双边理想等于 B。
\item
  \(k\)-代数 A 同构于 \(e\mathbf{M}_n(\mathrm{B})e\)。
\end{enumerate}

若条件 (i) 与 (ii) 满足，则 \((e\mathbf{M}_n(\mathrm{B})e,\mathrm{B})_k\)-双模 \(e\mathrm{B}_d^n\) 是可逆的。

鉴于定理 1（VIII, 第 101 页），k-代数 A 与 B 森田等价当且仅当它同构于一个投射的、有限生成的且是生成模的右 B-模的自同态代数。于是该命题由下面的引理 3 与引理 4 得出。

引理 3. —— 一个右 B-模 P 是投射的、有限生成的且是生成模，当且仅当存在一个整数 \(n \geqslant 0\) 与一个幂等元 \(e = (e_{ij})\)（在 \(\mathbf{M}_n(\mathrm{B})\) 中），具有下列性质：

\begin{enumerate}
\def\labelenumi{(\roman{enumi})}
\item
  B-模 P 同构于 \(e\mathrm{B}_d^n\)。
\item
  由诸元素 \(e_{ij}\) 生成的 B 的双边理想等于 B。
\end{enumerate}

设 P 是一个右 B-模。那么 P 是投射的且有限生成的，当且仅当它同构于一个形如 \(B_{d}^{n}\) 的 B-模（n 是一个 \(\geqslant 0\) 的整数）的直因子子模（II, §2, No.~2, 第 232 页，推论 1）。若我们把 k-代数 \(\mathbf{M}_{n}(\mathrm{B})\) 与 \(\operatorname{End}(\mathrm{B}_{d}^{n})\) 等同起来，则这就意味着存在 \(\mathbf{M}_{n}(\mathrm{B})\) 中的一个幂等元 e，使 P 同构于 \(eB_{d}^{n}\)。

B-模 P 是生成模当且仅当它的迹理想 \(\tau(\mathrm{P})\) 等于 B，即 \(\tau(e\mathrm{B}_{d}^{n})=\mathrm{B}\)（VIII, 第 80 页，定理 1）。设 \(x_{1},\ldots,x_{n}\) 是 \(B_{d}^{n}\) 中对应于矩阵 e 的各列的元素，并设 \(x_{i}^{*}\)（对 \(1\leqslant i\leqslant n\)）是线性形式 \((b_{1},\ldots,b_{n})\mapsto b_{i}\)（在 \(eB_{d}^{n}\) 上）。族 \((x_{1},\ldots,x_{n})\) 生成 B-模 \(eB_{d}^{n}\)，而族 \((x_{1}^{*},\ldots,x_{n}^{*})\) 生成它的对偶。而我们由 \(\langle x_{i}^{*},x_{j}\rangle=e_{ij}\)，故 \(\tau(eB_{d}^{n})\) 是由诸 \(e_{ij}\) 生成的 B 的双边理想。这就证明了引理 3。

引理 4. —— 设 V 是一个 B-模，E 是 V 的自同态 k-代数。设 e 是 V 的一个投影算子，P 是 e 的像。把 \(v \in eEe\) 映到 B-模 P 的自同态 \(x \mapsto v(x)\) 的映射是从 eEe 到 \(\operatorname{End}_{\mathrm{B}}(\mathrm{P})\) 的 k-代数同构。

用 \(\varphi:eEe\to\operatorname{End}_{\mathrm{B}}(\mathrm{P})\) 表示陈述中所描述的映射；它是一个 k-代数同态。设 \(u\in\operatorname{End}_{\mathrm{B}}(\mathrm{P})\)。用 v 表示由 \(v(x)=u(e(x))\)（\(x\in V\)）定义的 V 的自同态。我们有 \((eve)(x)=u(x)\)（对 \(x\in P\)），即 \(\varphi(eve)=u\)。于是 \(\varphi\) 是满射。设 w 是 \(\varphi\) 的核的一个元素；w 在 e 的核上与像上的限制均为零，故 w 为零，这就证明 \(\varphi\) 是单射。

例. —— 1) 设 A 是一个 k-代数，e 是 A 中满足 AeA = A 的幂等元。k-代数 eAe 可以等同于子模 \(eA_{d}\)（\(A_{d}\) 的）的自同态 k-代数（引理 4）。由于 AeA = A，由命题 14 推出 \(eA_{d}\) 是一个可逆 \((eAe, A)_{k}\)-双模，于是 \(k\)-代数 eAe 与 A 是森田等价的。此外，森田定理（VIII, 第 103 页）蕴含下面的结果：

\begin{enumerate}
\def\labelenumi{\alph{enumi})}
\item
  设 M 与 N 是左 A-模。每个从 eM 到 eN 的 eAe-线性映射都唯一地扩张为一个从 M 到 N 的 A-线性映射。
\item
  k-代数 eAe 上的每个左模都同构于形如 eM 的模，其中 M 是一个左 A-模。
\end{enumerate}

\begin{enumerate}
\def\labelenumi{\arabic{enumi})}
\setcounter{enumi}{1}
\tightlist
\item
  设 A 是一个 k-代数，\(n \geqslant 1\) 是一个整数。我们把矩阵代数 \(\mathbf{M}_{n}(\mathrm{A})\) 等同于右 A-模 \(A_{d}^{n}\) 的自同态代数。我们已经看到 A 与 \(\mathbf{M}_{n}(\mathrm{A})\) 是森田等价的。对每个左 A-模 M，把 \(A_{d}^{n} \otimes_{A} M\) 与 \(M^{n}\) 等同起来。于是代数 \(\mathbf{M}_{n}(\mathrm{A})\) 在 \(M^{n}\) 上有一个左作用，并且我们有
\end{enumerate}

\[
(a \cdot m) _ {i} = \sum_ {j = 1} ^ {n} a _ {i j} m _ {j}
\]

其中 \(a = (a_{ij})\) 属于 \(\mathbf{M}_n(\mathrm{A})\)，\(m = (m_i)\) 属于 \(\mathrm{M}^n\)。森田定理蕴含下面的结果：

\begin{enumerate}
\def\labelenumi{\alph{enumi})}
\item
  代数 \(\mathbf{M}_n(\mathrm{A})\) 上的每个左模都同构于形如 \(\mathrm{M}^n\) 的模，其中 \(\mathrm{M}\) 是一个左 A-模。
\item
  设 \(\mathrm{M}\) 是一个左 A-模。映射 \(\mathrm{N} \mapsto \mathrm{N}^n\) 是从 \(\mathrm{M}\) 的 A-子模之集到 \(\mathbf{M}_n(\mathrm{A})\)-子模之集（\(\mathrm{M}^n\) 的）的双射。
\item
  设 M 与 N 是左 A-模。对 A-线性映射 \(g: \mathrm{M} \to \mathrm{N}\)，用 \(g_{n}\) 表示映射 \((m_{i}) \mapsto (g(m_{i}))\)（从 \(\mathrm{M}^{n}\) 到 \(\mathrm{N}^{n}\)）。那么映射 \(g \mapsto g_{n}\) 是从 \(\operatorname{Hom}_{\mathrm{A}}(\mathrm{M}, \mathrm{N})\) 到 \(\operatorname{Hom}_{\mathbf{M}_{n}(\mathrm{A})}(\mathrm{M}^{n}, \mathrm{N}^{n})\) 的双射。
\item
  设 M 是一个左 A-模。模 \(M^{n}\)（在环 \(\mathbf{M}_{n}(\mathrm{A})\) 上）是不可分解的（相应地，半单的、单的、阿廷的、诺特的、有限生成的）当且仅当 A-模 M 亦然。
\end{enumerate}

\begin{enumerate}
\def\labelenumi{\arabic{enumi})}
\setcounter{enumi}{2}
\tightlist
\item
  设 A 是一个主理想整环，L 是一个非零的、有限生成的自由 A-模。设 B 是 L 的自同态环；则 L 是一个可逆 \((\mathrm{A},\mathrm{B})_{\mathbf{Z}}\)-双模，且环 A 与 B 是森田等价的。由森田定理、命题 10, a)（VIII, 第 109 页）以及有限生成 A-模的结构定理（VII, §4, No.~4, 第 19 页，定理 2），每个有限生成 B-模都同构于 \(\oplus_{i=1}^{m}(\mathrm{L}/\mathfrak{a}_{i}\mathrm{L})\)，其中 m 是一个自然数，诸 \(a_{i}\) 是 A 的理想，满足 \(a_{1} \subset a_{2} \subset \cdots \subset a_{m}\) 且 \(a_{m} \neq A\)；整数 m 与诸理想 \(a_{i}\) 是唯一确定的。由 VIII, 第 105 页的命题 6，B 的每个双边理想都形如 dB，其中 d 是 A 的一个元素。
\end{enumerate}

